\subsection{从可观察现象反推连接错误}

具体机器模型还应支持诊断。若上电后没有得到结果，不能只说“机器没跑起来”，而要利用
PC、事务和状态寄存器把故障范围逐步缩小。

\subsection{复位后连第一笔请求都没有}

若示波观察不到 \code{req\_valid}，先检查时钟是否产生边沿、复位是否释放以及 PC 是否
得到规定值。此时还没有理由检查任务映像，因为固件第一条取指尚未发生。

\subsection{请求地址正确，但始终没有响应}

观察到地址 \code{0x00000000} 和读请求，却没有 \code{resp\_valid}，故障位于互连选择、
启动存储连接或目标握手。若互连错误地把该地址送往 RAM，RAM 中的随机字节可能返回成功，
随后表现为无效操作码。区分“无响应”和“错误目标响应”需要同时观察目标选择信号。

\subsection{固件运行，但接收状态永远为零}

若 PC 已在固件轮询循环，\code{STATUS} 读也持续成功，问题不在处理器取指和 RAM。应检查
外部发送端是否产生串行位、控制器采样时序是否匹配以及接收移位寄存器是否能置
\code{RX\_READY}。软件看到的零值只是控制器没有宣布完整字节，不足以判断线路哪一端
失效。

\subsection{任务结果正确，外部端却没有收到}

先读 \code{MEM32[0x00102020]}。若值为 21，计算和 RAM 写入已经完成；再看 PC 是否到达
固件返回入口，区分“任务未返回”和“返回后发送失败”；最后检查
\code{TX\_READY/TX\_EMPTY} 与远端帧校验。诊断顺序沿完成边界向外推进。

\begin{longtable}{|L{0.25\linewidth}|L{0.27\linewidth}|L{0.34\linewidth}|}
\hline
最后确认成功的事实 & 下一项观察 & 尚不应检查的远端事项 \\ \hline
时钟与复位正常 & 第一笔取指请求 & 任务求和结果 \\ \hline
固件取指成功 & 串行状态与接收字节 & 任务返回地址 \\ \hline
映像校验成功 & 初始现场与入口 PC & 外部结果帧 \\ \hline
结果槽为 21 & 返回入口 PC & 传输校验 \\ \hline
\code{TX\_EMPTY=1} & 远端接收与帧校验 & RAM 计算过程 \\ \hline
\end{longtable}

这种从已知成功边界向下一边界推进的方法，与后续操作系统中的系统调用、设备请求和文件
持久化诊断具有相同结构：先确定事实在哪一层成立，再检查下一层所需的新状态。

\subsection{最小机器的最终接线清单}

到本章末再列清单，作用是复核已经推导的连接，而不是用名词代替正文：

\begin{longtable}{|L{0.21\linewidth}|L{0.31\linewidth}|L{0.34\linewidth}|}
\hline
连接 & 承载的信息 & 接错后的直接后果 \\ \hline
时钟到处理器、互连和控制器 & 共同状态更新边沿 & 握手双方不能在同一边界观察状态 \\ \hline
复位到处理器与控制状态 & PC 起点、协议初态 & 旧事务或任意 PC 被当作有效 \\ \hline
处理器请求到互连 & 地址、读写、宽度、写数据 & 无法取指或访问错误目标 \\ \hline
目标响应经互连回处理器 & 状态与读数据 & 指令永久等待或读取错误字节 \\ \hline
互连到启动存储 & 固件读请求 & 复位后没有可靠指令 \\ \hline
互连到 RAM & 任务取指和数据读写 & 任务无法装入或保存运行状态 \\ \hline
互连到串行控制器 & 状态、收发字节 & 外部任务无法进入或结果无法报告 \\ \hline
串行控制器到外部端 & 按时间传送的位 & 软件寄存器与外部信息断开 \\ \hline
\end{longtable}

每一行都能回到前文的一次具体事务。清单中没有中断控制器、计时器、地址转换器、磁盘或
操作系统；这些结构只有在后续问题确实需要时才会加入，并重新说明硬件模型、软件接口和
提供的抽象。
